\section{Sources}\label{app:sources}

We compared the bibliographic data of every reference with the publisher's record, with Crossref
or with arXiv, except where we say otherwise. For~\cite{ES} we used the scan that is linked from
the OEIS. We read~\cite{AACST}, \cite{ES}, \cite{Griffin}, \cite{KNT}, \cite{Markstrom}, and the
parts of~\cite{HK} that concern cycles. Of~\cite{AT} we read the abstract. We have not
seen~\cite{Bass}, \cite{Bondy}, \cite{EP}, \cite{Guichard} or~\cite{KS}. What we say
about~\cite{Bondy} is taken from~\cite{Bloom} and~\cite{KNT}. What we say about~\cite{Guichard} is
taken from~\cite{Markstrom}, and we did not compare its bibliographic data with a record. No proof
in this paper depends on a source that we have not seen. The lemma from~\cite{EP} is proved in
Section~\ref{sec:exact}, and the facts related to~\cite{Bass,KS} are proved in
Section~\ref{sec:forest}.

We looked for an earlier proof of $\psi(r)\sim2^{r-1}$, for any bound below $\frac{15}{16}2^{r}$
for cubic graphs, and for an upper bound of the form $2^{r-1}+o(2^{r})$ with an explicit error
term. We used web searches and the literature discussed in~\cite{AACST}, \cite{HK}
and~\cite{Markstrom}, and we found none. A search of this kind can miss things.

The two theorems of Section~\ref{sec:exact} are elementary, and they are close to known results.
Lemma~\ref{lem:EP} is due to Erd\H{o}s and P\'osa. The graphs in which every nonempty even subgraph
is a cycle were characterised by Maurer and by Mateti and Deo; we know this
from~\cite[Theorem~1]{ES} and have not seen the original papers. We do not know whether
Proposition~\ref{prop:terminal} or the bound in Proposition~\ref{prop:nowaste} has been stated
before.

\paragraph{The supplement.}
The entries~\cite{S0} to~\cite{S5} are files in the folder
\lean{paper/cycle-count/research/2026-11-10} of the repository~\cite{Repo}. They are research
reports and notes, written by the two models during the work. They are not publications, and no
person has reviewed them. The reports~\cite{S0} to~\cite{S4} have compiled versions with the same
name and the ending \lean{.pdf}. The same folder contains the other notes of this work, among them
results that this paper does not use. The repository also keeps the earlier drafts of this paper,
next to the manuscript in the folder \lean{paper/cycle-count}.

\begin{thebibliography}{99}\small

\bibitem{AACST} A.~Alahmadi, R.~E.~L.~Aldred, R.~dela~Cruz, P.~Sol\'e and C.~Thomassen, The maximum
number of minimal codewords in long codes, \emph{Discrete Appl.\ Math.}\ \textbf{161} (2013),
424--429. arXiv:1201.6043.

\bibitem{AT} R.~E.~L.~Aldred and C.~Thomassen, On the maximum number of cycles in a planar graph,
\emph{J.~Graph Theory} \textbf{57} (2008), 255--264.

\bibitem{Bass} H.~Bass, The Ihara--Selberg zeta function of a tree lattice, \emph{Internat.\
J.~Math.}\ \textbf{3} (1992), 717--797.

\bibitem{Bloom} T.~F.~Bloom, Erd\H{o}s Problem \#1016, \url{https://www.erdosproblems.com/1016},
accessed November 2026.

\bibitem{Bondy} J.~A.~Bondy, Pancyclic graphs~I, \emph{J.~Combin.\ Theory Ser.~B} \textbf{11}
(1971), 80--84.

\bibitem{ES} R.~C.~Entringer and P.~J.~Slater, On the maximum number of cycles in a graph,
\emph{Ars Combin.}\ \textbf{11} (1981), 289--294. A scan is linked from the entry A006188 of the
OEIS, \url{https://oeis.org/A006188}.

\bibitem{EP} P.~Erd\H{o}s and L.~P\'osa, On the maximal number of disjoint circuits of a graph,
\emph{Publ.\ Math.\ Debrecen} \textbf{9} (1962), 3--12.

\bibitem{Griffin} S.~Griffin, Minimal pancyclicity, arXiv:1312.0274 (2013).

\bibitem{Guichard} D.~R.~Guichard, The maximum number of cycles in graphs, \emph{Congr.\ Numer.}\
\textbf{121} (1996), 211--215.

\bibitem{HK} P.~Horak and D.~Kim, Connected cubic graphs with the maximum number of perfect
matchings, \emph{J.~Graph Theory} \textbf{99} (2022), 671--690. arXiv:2006.13459.

\bibitem{KNT} KNT, The minimum number of edges in a pancyclic graph, manuscript, September 2026,
\url{https://jehlp.net/erdos1016/}.

\bibitem{Claim} KNT, Proof claim for Problem \#1016, Erd\H{o}s Problems forum, 24 September 2026,
\url{https://www.erdosproblems.com/forum/thread/1016/proof-claims}, accessed November 2026.

\bibitem{Repo} KNT, Astra and Fable~5.2, Lean~4 formalisation of the results of~\cite{KNT} and of
parts of this paper, \url{https://github.com/JWKNT/erdos1016}.

\bibitem{KS} M.~Kotani and T.~Sunada, Zeta functions of finite graphs, \emph{J.~Math.\ Sci.\ Univ.\
Tokyo} \textbf{7} (2000), 7--25.

\bibitem{Lean4} L.~de~Moura and S.~Ullrich, The Lean~4 theorem prover and programming language,
in: \emph{CADE~28}, Lecture Notes in Comput.\ Sci.\ \textbf{12699}, Springer, 2021, 625--635.

\bibitem{Markstrom} K.~Markstr\"om, Extremal graphs for some problems on cycles in graphs,
\emph{Congr.\ Numer.}\ \textbf{171} (2004), 179--192. We read the author's version of August 2005.

\bibitem{Mathlib} The mathlib Community, The Lean mathematical library, in: \emph{CPP~2020}, ACM,
2020, 367--381.

\bibitem{S0} Astra, The unrestricted sharp-order cycle bound, research report, 11 November 2026.
File \lean{unrestricted-sharp-cycle-bound.tex} in the supplement. The proof that
Sections~\ref{sec:sharp} to~\ref{sec:halfproof} follow.

\bibitem{S1} Astra, The sharp second-order cycle bound for cyclic cuts of size at most five,
research report, 10 November 2026. File \lean{small-cut-sharp-bound.tex} in the supplement. Proof
of Theorem~\ref{thm:smallcut}.

\bibitem{S2} Astra, A strict cycle deficit at bounded girth, research report, 11 November 2026.
File \lean{bounded-girth-cycle-deficit.tex} in the supplement. Proof of Theorem~\ref{thm:girth}.

\bibitem{S3} Astra and Fable~5.2, Additive stability for terminal path counts, research report,
10 November 2026. File \lean{terminal-path-stability.tex} in the supplement. Proof of
Theorem~\ref{thm:stability}.

\bibitem{S4} Astra, A polynomial error for the number of cycles, research report, 10 November
2026, file \lean{polynomial-cycle-error.tex}; and Fable~5.2, Link components are independent,
research note, file \lean{fable-parallel-links.md}. Both in the supplement.

\bibitem{S5} Astra and Fable~5.2, research notes, November 2026. Files \lean{no-waste-phase2.tex}
(the constant in Proposition~\ref{prop:nowaste}) and \lean{fable-census.md} (cubic graphs with at
most $22$ vertices) in the supplement.

\bibitem{S6} Astra, The eventual exact maximum number of cycles, research report, 11 November
2026. File \lean{eventual-exact-cycle-maximum.tex} in the folder
\lean{paper/cycle-count/research/2026-11-11} of the repository~\cite{Repo}. Not used in this paper.

\end{thebibliography}
